\ifdefined\gkver\else\def\gkver{student}\fi
\documentclass[\gkver]{gaokaozhenti}
\newcommand{\ccnum}[1]{{\fontspec{Noto Sans CJK JP}#1}}
\begin{document}

\chapter{2013年福建理综}
%% paperType: 理综物理部分

\begin{enumerate}
\item
%% number: 13
%% paperName: 2013年福建理综第 13 题 6 分
%% typeId: 1
%% type: 单选题
%% chapter: 曲线运动
%% point: 天体运动
%% method: 无
%% score: 6
%% degree: 650
%% duplicateId: 0
%% body:
设太阳质量为 $M$，某行星绕太阳公转周期为 $T$，轨道可视作半径为 $r$ 的圆。已知万有引力常量为 $G$，则描述该行星运动的上述物理量满足\xzanswer{A}
\fourchoices[answer=A]
{$GM=\frac{4\pi^{2}r^{3}}{T^{2}}$}
{$GM=\frac{4\pi^{2}r^{2}}{T^{2}}$}
{$GM=\frac{4\pi^{2}r^{2}}{T^{3}}$}
{$GM=\frac{4\pi r^{3}}{T^{2}}$}
%% answer: A
%% memo:
\memoanswer{
本题原卷及材料中未提供详解，待补充。
}

\item
%% number: 14
%% paperName: 2013年福建理综第 14 题 6 分
%% typeId: 1
%% type: 单选题
%% chapter: 光学
%% point: 光的折射
%% method: 无
%% score: 6
%% degree: 650
%% duplicateId: 0
%% body:
一束由红、紫两色光组成的复色光，从空气斜射向玻璃三棱镜。下面四幅图中能正确表示该复色光经三棱镜分离成两束单色光的是\xzanswer{B}
\fourchoices[answer=B, ispicture=true, h=1.6cm]
{figs/14a.png}
{figs/14b.png}
{figs/14c.png}
{figs/14d.png}
%% answer: B
%% memo:
\memoanswer{
本题原卷及材料中未提供详解，待补充。
}

\item
%% number: 15
%% paperName: 2013年福建理综第 15 题 6 分
%% typeId: 1
%% type: 单选题
%% chapter: 交流电
%% point: 交流电
%% method: 无
%% score: 6
%% degree: 650
%% duplicateId: 0
%% body:
如图，实验室一台手摇交流发电机，内阻 $r=1.0\UO$，外接 $R=9.0\UO$ 的电阻。闭合开关 S，当发电机转子以某一转速匀速转动时，产生的电动势 $e=10\sqrt{2}\sin 10\pi t$（$\UV$），则\xzanswer{D}

\onepicture[w=3cm, align=right]{figs/15.png}

\fourchoices[answer=D]
{该交变电流的频率为 $10\UHz$}
{该电动势的有效值为 $10\sqrt{2}\UV$}
{外接电阻 $R$ 所消耗的电功率为 $10\UW$}
{电路中理想交流电流表 A 的示数为 $1.0\UA$}
%% answer: D
%% memo:
\memoanswer{
本题原卷及材料中未提供详解，待补充。
}

\item
%% number: 16
%% paperName: 2013年福建理综第 16 题 6 分
%% typeId: 1
%% type: 单选题
%% chapter: 机械波
%% point: 波的图象
%% method: 无
%% score: 6
%% degree: 650
%% duplicateId: 0
%% body:
如图，$t=0$ 时刻，波源在坐标原点从平衡位置沿 $y$ 轴正方向开始振动，振动周期为 $0.4\Us$，在同一均匀介质中形成沿 $x$ 轴正、负两方向传播的简谐横波。下图中能够正确表示 $t=0.6\Us$ 时波形的图是\xzanswer{C}
\fourchoices[answer=C, ispicture=true, h=1.6cm]
{figs/16a.png}
{figs/16b.png}
{figs/16c.png}
{figs/16d.png}
%% answer: C
%% memo:
\memoanswer{
本题原卷及材料中未提供详解，待补充。
}

\item
%% number: 17
%% paperName: 2013年福建理综第 17 题 6 分
%% typeId: 1
%% type: 单选题
%% chapter: 运动和力
%% point: 力学单位制
%% method: 无
%% score: 6
%% degree: 850
%% duplicateId: 0
%% body:
在国际单位制（简称 SI）中，力学和电学的基本单位有：m（米）、kg（千克）、s（秒）、A（安培）。导出单位 V（伏特）用上述基本单位可表示为\xzanswer{B}
\fourchoices[answer=B]
{m$^{2}\cdot$kg$\cdot$s$^{-4}\cdot$A$^{-1}$}
{m$^{2}\cdot$kg$\cdot$s$^{-3}\cdot$A$^{-1}$}
{m$^{2}\cdot$kg$\cdot$s$^{-2}\cdot$A$^{-1}$}
{m$^{2}\cdot$kg$\cdot$s$^{-1}\cdot$A$^{-1}$}
%% answer: B
%% memo:
\memoanswer{
本题原卷及材料中未提供详解，待补充。
}

\item
%% number: 18
%% paperName: 2013年福建理综第 18 题 6 分
%% typeId: 1
%% type: 单选题
%% chapter: 电磁感应
%% point: 线圈过有界磁场
%% method: 无
%% score: 6
%% degree: 450
%% duplicateId: 0
%% body:
如图，矩形闭合线框在匀强磁场上方，由不同高度静止释放，用 $t_{1}$、$t_{2}$ 分别表示线框 $ab$ 边和 $cd$ 边刚进入磁场的时刻。线框下落过程形状不变，$ab$ 边始终保持与磁场水平边界 $OO'$ 平行，线框平面与磁场方向垂直。设 $OO'$ 下方磁场区域足够大，不计空气影响，则下列哪一个图像不可能反映线框下落过程中速度 $v$ 随时间 $t$ 变化的规律\xzanswer{A}

\onepicture[w=3cm, align=right]{figs/18.png}

\fourchoices[answer=A, ispicture=true, h=1.8cm]
{figs/18a.png}
{figs/18b.png}
{figs/18c.png}
{figs/18d.png}
%% answer: A
%% memo:
\memoanswer{
本题原卷及材料中未提供详解，待补充。
}

\item
%% number: 19
%% paperName: 2013年福建理综第 19 题 12 分
%% typeId: 4
%% type: 实验题
%% chapter: 电路
%% point: 电路实验
%% method: 无
%% score: 12
%% degree: 450
%% duplicateId: 0
%% body:
\begin{enumerate}
	\item
	（6 分）在“探究恒力做功与动能改变的关系”实验中（装置如图甲）：
	\begin{enumerate}
		\item
		下列说法哪一项是正确的\tkanswer[1cm]{C}。（填选项前字母）

		（A）平衡摩擦力时必须将钩码通过细线挂在小车上

		（B）为减小系统误差，应使钩码质量远大于小车质量

		（C）实验时，应使小车靠近打点计时器由静止释放
		\item
		图乙是实验中获得的一条纸带的一部分，选取 $O$、$A$、$B$、$C$ 计数点，已知打点计时器使用的交流电频率为 $50\UHz$，则打 $B$ 点时小车的瞬时速度大小为\tkanswer[1.4cm]{$0.653$}$\Ums$（保留三位有效数字）。
	\end{enumerate}

	\onepicture[w=6cm, align=right]{figs/19.png}

	\item
	（12 分）硅光电池在无光照时不产生电能，可视为一电子元件。某实验小组设计如图甲电路，给硅光电池加反向电压（硅光电池负极接高电势点，正极接低电势点），探究其在无光照时的反向伏安特性。图中电压表 $\mathrm{V}_{1}$ 量程选用 $3\UV$，内阻为 $6.0\UkO$；电压表 $\mathrm{V}_{2}$ 量程选用 $15\UV$，内阻约为 $30\UkO$；$R_{0}$ 为保护电阻；直流电源电动势 $E$ 约为 $12\UV$，内阻不计。

	\onepicture[w=11cm]{figs/19a.png}

	\begin{enumerate}
		\item
		根据图甲，用笔画线代替导线，将图乙连接成完整电路。
		\item
		用遮光罩罩住硅光电池，闭合开关 $S$，调节变阻器 $R$，读出电压表 $\mathrm{V}_{1}$、$\mathrm{V}_{2}$ 的示数 $U_{1}$、$U_{2}$。
		\item
		（i）某次测量时，电压表 $\mathrm{V}_{1}$ 示数如图丙，则 $U_{1}=$\tkanswer[1.4cm]{$1.40$}$\UV$，可算出通过硅光电池的反向电流大小为\tkanswer[1.4cm]{$0.23$}$\UmA$（保留两位小数）。
		\item
		（ii）该小组测出大量数据，筛选出下表所示的 9 组 $U_{1}$、$U_{2}$ 数据，算出相应的硅光电池两端反向电压 $U_{x}$ 和通过反向电流 $I_{x}$（表中“－”表示反向），并在坐标纸上建立 $I_{x}-U_{x}$ 坐标系，标出了与表中前 5 组 $U_{x}$、$I_{x}$ 数据对应的 5 个坐标点。请你标出余下的 4 个坐标点，并绘出 $I_{x}-U_{x}$ 图线。
	\end{enumerate}

	\begin{table}[htbp]
	\centering
	\begin{tblr}{|c|c|c|c|c|c|c|c|c|c|}
	\hline
	  & 1 & 2 & 3 & 4 & 5 & 6 & 7 & 8 & 9 \\
	\hline
	$U_{1}/\UV$ & 0.00 & 0.00 & 0.06 & 0.12 & 0.24 & 0.42 & 0.72 & 1.14 & 1.74 \\
	\hline
	$U_{1}/\UV$ & 0.0 & 1.0 & 2.1 & 3.1 & 4.2 & 5.4 & 6.7 & 8.1 & 9.7 \\
	\hline
	$U_{x}/\UV$ & 0.0 & -1.0 & -2.0 & -3.0 & -4.0 & -5.0 & -6.0 & -7.0 & -8.0 \\
	\hline
	$I_{x}/\UmA$ & -0.00 & -0.00 & -0.01 & -0.02 & -0.04 & -0.07 & -0.12 & -0.19 & -0.29 \\
	\hline
	\end{tblr}
	\end{table}

	\onepicture[w=5.5cm]{figs/19b.png}

	（iii）由 $I_{x}-U_{x}$ 图线可知，硅光电池无光照下加反向电压时，$I_{x}$ 与 $U_{x}$ 成\tkanswer[1.2cm]{非线性}（填“线性”或“非线性”）关系。
\end{enumerate}
%% answer: 
\jdanswer{
（1）① C；② $0.653\Ums$。

（2）①电路连接如图所示；②（i）$U_{1}=1.40\UV$，反向电流大小为 $0.23\UmA$；（ii）描点并作图如图所示；（iii）非线性。
}
%% memo:
\memoanswer{
本题原卷及材料中未提供详解，待补充。
}

\item
%% number: 20
%% paperName: 2013年福建理综第 20 题 15 分
%% typeId: 6
%% type: 计算题
%% chapter: 功和能
%% point: 机械能守恒定律
%% method: 无
%% score: 15
%% degree: 650
%% duplicateId: 0
%% body:
如图，一不可伸长的轻绳上端悬挂于 $O$ 点，下端系一质量 $m=1.0\Ukg$ 的小球。现将小球拉到 $A$ 点（保持绳绷直）由静止释放，当它经过 $B$ 点时绳恰好被拉断，小球平抛后落在水平地面上的 $C$ 点。地面上的 $D$ 点与 $OB$ 在同一竖直线上，已知绳长 $L=1.0\Um$，$B$ 点离地高度 $H=1.0\Um$，$A$、$B$ 两点的高度差 $h=0.5\Um$，重力加速度 $g$ 取 $10\Umsq$，不计空气影响，求：
\begin{enumerate}
	\item
	地面上 $DC$ 两点间的距离 $s$；
	\item
	轻绳所受的最大拉力大小。
\end{enumerate}

\onepicture[w=5cm, align=right]{figs/20.png}

%% answer: 
\jdanswer{
\begin{enumerate}
	\item
	$s=1.41\Um$

	\item
	$F=20\UN$
\end{enumerate}
}
%% memo:
\memoanswer{
（1）小球从 A 到 B 过程机械能守恒，有

$mgh=\frac{1}{2}mv_{B}^{2}$\ccnum{①}，

小球从 B 到 C 做平抛运动，在竖直方向上有

$H=\frac{1}{2}gt^{2}$\ccnum{②}，

在水平方向上有

$s=v_{B}t$\ccnum{③}，

由①②③式解得 $s=1.41\Um$\ccnum{④}。

（2）小球下摆到达 B 点时，绳的拉力和重力的合力提供向心力，由①有

$F-mg=m\frac{v_{B}^{2}}{L}$\ccnum{⑤}，

由①⑤式解得 $F=20\UN$。

根据牛顿第三定律 $F'=-F$，轻绳所受的最大拉力为 $20\UN$。
}

\item
%% number: 21
%% paperName: 2013年福建理综第 21 题 19 分
%% typeId: 1
%% type: 单选题
%% chapter: 运动和力
%% point: 牛顿第二定律
%% method: 无
%% score: 19
%% degree: 850
%% duplicateId: 0
%% body:
质量为 $M$、长为 $\sqrt{3}L$ 的杆水平放置，杆两端 A、B 系着长为 $3L$ 的不可伸长且光滑的柔软轻绳，绳上套着一质量为 $m$ 的小铁环。已知重力加速度为 $g$，不计空气影响。
\begin{enumerate}
	\item
	现让杆和环均静止悬挂在空中，如图甲，求绳中拉力的大小；
	\item
	若杆与环保持相对静止，在空中沿 $AB$ 方向水平向右做匀加速直线运动，此时环恰好悬于 $A$ 端的正下方，如图乙所示。
	\begin{enumerate}
		\item
		求此状态下杆的加速度大小 $a$；
		\item
		为保持这种状态需在杆上施加一个多大的外力，方向如何？
	\end{enumerate}
\end{enumerate}

\onepicture[w=5cm, align=right]{figs/21.png}

%% answer: 
\jdanswer{
\begin{enumerate}
	\item
	$F_{T}=\frac{\sqrt{6}}{4}mg$

	\item
	① $a=\frac{\sqrt{3}}{3}g$；② $F=\frac{2\sqrt{3}}{3}(m+M)g$，方向与水平方向夹角为 $60^{\circ}$ 斜向右上方。
\end{enumerate}
}
%% memo:
\memoanswer{
（1）如图 1，设平衡时，绳中拉力为 $F_{T}$，有

$2F_{T}\cos\theta-mg=0$\ccnum{①}，

\onepicture[w=4cm]{figs/21_an1.jpg}

由图知 $\cos\theta=\frac{\sqrt{6}}{3}$\ccnum{②}，

由①②式解得 $F_{T}=\frac{\sqrt{6}}{4}mg$\ccnum{③}。

（2）①此时，对小铁环受力分析如图 2，有

\onepicture[w=4cm]{figs/21_an2.jpg}

$F_{T}''\sin\theta'=ma$\ccnum{④}，

$F_{T}'+F_{T}''\cos\theta'-mg=0$\ccnum{⑤}，

由图知 $\theta'=60^{\circ}$，代入④⑤式解得 $a=\frac{\sqrt{3}}{3}g$\ccnum{⑥}。

②如图 3，设外力 $F$ 与水平方向成 $\alpha$ 角，将杆和小铁环当成一个整体，有

\onepicture[w=4cm]{figs/21_an3.jpg}

$F\cos\alpha=(M+m)a$\ccnum{⑦}，

$F\sin\alpha-(M+m)g=0$\ccnum{⑧}，

由⑥⑦⑧式解得 $F=\frac{2\sqrt{3}}{3}(m+M)g$，$\tan\alpha=\sqrt{3}$（或 $\alpha=60^{\circ}$），即与水平方向夹角为 $60^{\circ}$ 斜向右上方。
}

\item
%% number: 22
%% paperName: 2013年福建理综第 22 题 20 分
%% typeId: 6
%% type: 计算题
%% chapter: 磁场
%% point: 带电粒子在磁场中的运动
%% method: 无
%% score: 20
%% degree: 450
%% duplicateId: 0
%% body:
如图甲，空间存在一范围足够大的垂直于 $xOy$ 平面向外的匀强磁场，磁感应强度大小为 $B$。让质量为 $m$，电量为 $q$（$q<0$）的粒子从坐标原点 $O$ 沿 $xOy$ 平面以不同的初速度大小和方向入射到该磁场中。不计重力和粒子间的影响。
\begin{enumerate}
	\item
	若粒子以初速度 $v_{1}$ 沿 $y$ 轴正向入射，恰好能经过 $x$ 轴上的 $A$（$a$，0）点，求 $v_{1}$ 的大小；
	\item
	已知一粒子的初速度大小为 $v$（$v>v_{1}$），为使该粒子能经过 $A$（$a$，0）点，其入射角 $\theta$（粒子初速度与 $x$ 轴正向的夹角）有几个？并求出对应的 $\sin\theta$ 值；
	\item
	如图乙，若在此空间再加入沿 $y$ 轴正向、大小为 $E$ 的匀强电场，一粒子从 $O$ 点以初速度 $v_{0}$ 沿 $x$ 轴正向发射。研究表明：粒子在 $xOy$ 平面内做周期性运动，且在任一时刻，粒子速度的 $x$ 分量 $v_{x}$ 与其所在位置的 $y$ 坐标成正比，比例系数与场强大小 $E$ 无关。求该粒子运动过程中的最大速度值 $v_{m}$。
\end{enumerate}

\onepicture[w=10cm, align=right]{figs/22.png}

%% answer: 
\jdanswer{
\begin{enumerate}
	\item
	$v_{1}=\frac{Bqa}{2m}$

	\item
	$\sin\theta=\frac{aqB}{2mv}$

	\item
	$v_{m}=\frac{E}{B}+\sqrt{\frac{E^{2}}{B^{2}}+v_{0}^{2}}$
\end{enumerate}
}
%% memo:
\memoanswer{
（1）带电粒子以速率 $v$ 在匀强磁场 $B$ 中做匀速圆周运动，半径为 $R$，有

$qvB=m\frac{v^{2}}{R}$\ccnum{①}，

当粒子沿 $y$ 轴正向入射，转过半个圆周至 $A$ 点，该圆周的半径为 $R_{1}$，有

$R_{1}=\frac{a}{2}$\ccnum{②}，

由②代入①式得 $v_{1}=\frac{Bqa}{2m}$\ccnum{③}。

（2）如图，$O$、$A$ 两点处于同一圆周上，且圆心在 $x=\frac{a}{2}$ 的直线上，半径为 $R$。当确定一个初速率 $v$ 时，有 2 个入射角，分别在第 1、2 象限，有

\onepicture[w=5cm]{figs/22_an.jpg}

$\sin\theta'=\sin\theta=\frac{a}{2R}$\ccnum{④}，

由①④解得 $\sin\theta=\frac{aqB}{2mv}$\ccnum{⑤}。

（3）粒子在运动过程中仅有电场力做功，因此在轨道的最高点处速率最大，用 $y_{m}$ 表示其 $y$ 坐标，由动能定理，有

$qEy_{m}=\frac{1}{2}mv_{m}^{2}-\frac{1}{2}mv_{0}^{2}$\ccnum{⑥}，

由题意，有 $v_{m}=ky_{m}$\ccnum{⑦}，

若 $E=0$ 时，粒子以初速度 $v_{0}$ 沿 $y$ 轴正向入射，有

$qv_{0}B=m\frac{v_{0}^{2}}{R_{0}}$\ccnum{⑧}，

$v_{0}=kR_{0}$\ccnum{⑨}，

由⑥⑦⑧⑨式解得 $v_{m}=\frac{E}{B}+\sqrt{\frac{E^{2}}{B^{2}}+v_{0}^{2}}$。
}

\item
%% number: 29
%% paperName: 2013年福建理综第 29 题 6 分
%% typeId: 1
%% type: 单选题
%% chapter: 气体性质
%% point: 玻意耳定律
%% method: 无
%% score: 6
%% degree: 650
%% duplicateId: 0
%% body:
【物理——选修 3-3】（本题共有两小题，每小题 6 分，共 12 分。每小题只有一个选项符合题意）

（1）下列四幅图中，能正确反映分子间作用力 $f$ 和分子势能 $E_{p}$ 随分子间距离 $r$ 变化关系的图线是\tkanswer[1.2cm]{B}。（填选图下方的字母）

\fourchoices[ispicture=true, h=2.2cm]
{figs/29a.png}
{figs/29b.png}
{figs/29c.png}
{figs/29d.png}

（2）某自行车轮胎的容积为 $V$，里面已有压强为 $p_{0}$ 的空气，现在要使轮胎内的气压增大到 $p$，设充气过程为等温过程，空气可看作理想气体，轮胎容积保持不变，则还要向轮胎充入温度相同，压强也是 $p_{0}$，体积为\xzanswer{C} 的空气。（填选项前的字母）
\fourchoices[answer=C]
{$\frac{p_{0}}{p}V$}
{$\frac{p}{p_{0}}V$}
{（$\frac{p}{p_{0}}-1$）$V$}
{（$\frac{p}{p_{0}}+1$）$V$}
%% answer: C
%% memo:
\memoanswer{
本题原卷及材料中未提供详解，待补充。
}

\item
%% number: 30
%% paperName: 2013年福建理综第 30 题 6 分
%% typeId: 1
%% type: 单选题
%% chapter: 运动和力
%% point: 动量守恒定律
%% method: 无
%% score: 6
%% degree: 650
%% duplicateId: 0
%% body:
【物理——选修 3-5】（本题共有两小题，每小题 6 分，共 12 分。每小题只有一个选项符合题意）

（1）在卢瑟福 $\alpha$ 粒子散射实验中，金箔中的原子核可以看作静止不动，下列各图画出的是其中两个 $\alpha$ 粒子经历金箔散射过程的径迹，其中正确的是\tkanswer[1.2cm]{C}。（填选图下方的字母）

\fourchoices[ispicture=true, h=1.8cm]
{figs/30a.png}
{figs/30b.png}
{figs/30c.png}
{figs/30d.png}

（2）将静置在地面上，质量为 $M$（含燃料）的火箭模型点火升空，在极短时间内以相对地面的速度 $v_{0}$ 竖直向下喷出质量为 $m$ 的炽热气体。忽略喷气过程重力和空气阻力的影响，则喷气结束时火箭模型获得的速度大小是\xzanswer{D}。（填选项前的字母）
\fourchoices[answer=D]
{$\frac{m}{M}v_{0}$}
{$\frac{M}{m}v_{0}$}
{$\frac{M}{M-m}v_{0}$}
{$\frac{m}{M-m}v_{0}$}
%% answer: D
%% memo:
\memoanswer{
本题原卷及材料中未提供详解，待补充。
}

\end{enumerate}

\end{document}
